\section{Design Implications and Limits}
The mechanism yields a useful sequence of questions. Can a material perturbation enter the proposed current direction? Can that direction return to the receiver after retained-space scattering? Does the resulting finite material displacement reduce the present full quadratic? These questions concern different operators and cannot be replaced by one unconditional current ranking.

Conditional exchange starts from a stable receiver-derived space and changes its directions through a small number of checked replacements. No-op avoids accepting an unfavorable checked action; the substantive result is the measured decrease in absolute and object-summed full-GN gaps. A fixed current dimension controls representation size, not construction or verification cost. Each improvement purchases additional physics actions, whose setup, proposal, rejection, and offline evaluation costs remain visible.

The incoming pool restricts which opportunities can be tested. Full-dictionary reduced ranking can be accurate while the 12-atom pool omits a valuable direction, and successive exchanges can enter different local neighborhoods. Candidate coverage, scoring error, and path interactions are reported separately. Improving candidate acquisition is future work; the present method and validation do not add a learned proposer.

The GN metric measures fidelity to a specified regularized local update. It is not truth error, and a full GN step can itself be a poor estimator under limited data or inappropriate regularization. Constraints and nonlinear acceptance act on a different executed displacement. Their common line search is essential to interpreting transfer to a final material estimate. Likewise, an FP64 numerical significance threshold does not establish a deterministic output-error radius or a discretization-independent safety guarantee.

The Maxwell contribution is the explicit entry--feedback--return path of a specified current-space modification and its residual-dependent material utility. TSOM already incorporates internal propagation, and generic goal-oriented reduction already uses task information. The physical factorization, concrete dark-current and interference witnesses, and fixed-dimension exchange evidence provide the electromagnetic design content. Primary-source gaps preclude historical priority claims.

\section{Conclusion}
A receiver current space should be modified according to the material step that a specified replacement actually produces. Eliminating retained currents exposes the feedback-dressed observation and material injection factors, while the normal equations expose both curvature change and residual pullback. Their finite signed gain supplies a direct test on a common material quadratic. At unchanged endpoint current rank, the locked rule reduces the median summed gap by 50.7\% across ten complete additional Maxwell objects. The effect persists in the bounded noisy-residual tests and has a limited reconstruction consequence in three completed constrained pairs. The resulting design principle is operational: trace a current replacement through material excitation and receiver return, then check the finite material displacement before committing to the modified space.
